A central challenge in developing alternative materials is the narrow domain specialization of human scientists. Polymer scientists, for instance, are highly trained to optimize polymers within well-defined design frameworks, yet this very depth often comes at the cost of breadth. Without a cross-disciplinary perspective that bridges polymers with other materials, such as ceramics, metals, biomaterials, or emerging hybrid systems, experts may struggle to envision replacements that draw on principles or insights outside their immediate field. 

This siloed structure of human knowledge limits scientific discovery in two key ways. It narrows the ability to scan the full landscape of materials science and engineering for viable substitutes, and it reduces the chance encounters between disciplines where novel ideas often emerge. As a result, even when the demand for sustainable alternatives is urgent, the process of hypothesis generation remains slowed by the inherent bounds of human expertise and imagination.

The challenge is particularly visible in efforts to identify substitutes for materials that are difficult to replace. Among these are per- and polyfluoroalkyl substances, commonly referred to as PFAS, a class of synthetic chemicals valued for their durability, chemical resistance, and water- and oil-repellent properties \cite{shabanian_durable_2023, leung_physicochemical_2023}. These same characteristics contribute to the environmental persistence of certain PFAS compounds and have prompted increasing regulatory attention in parts of the world \cite{zheng_per-_2021, henderson_analysis_2024, cousins_high_2020}. Despite this evolving landscape, PFAS remain widely used across industries, as few alternative materials offer a comparable combination of properties.

The growing performance of large language models (LLMs) \cite{Vaswani2017attention, touvron_llama_2023, wei_emergent_2022, zhao_survey_2023, teubner_welcome_2023, naveed_comprehensive_2023} presents a compelling strategy to reason over the massive corpus of scientific knowledge \cite{zhang_exploring_2025, rueda_understanding_2025} in ways that extend beyond the limitations of human-driven analysis. In particular, LLM-driven multi-agent frameworks enable the orchestration of domain-specialized agents to interact and coordinate toward reaching a common goal. These systems are poised for the requirements of materials design problems, which fundamentally demand a multi-step workflow wherein agents work together like a group of collaborative scientists. Unlike traditional machine learning techniques where a single AI model trains once on relevant data, these multi-agent systems can adapt to incoming data and solicit new data on the fly to perform complex tasks. X-LoRA demonstrates how LLMs can be continually adapted in a modular, mixture-of-experts manner by routing each query through specialized low-rank adapter ``experts'', enabling efficient, on-the-fly domain specialization (e.g., for protein mechanics and molecular design) without retraining the full base model~\cite{Buehler2024XLoRA}. Agents can also be integrated with external tools such as physics-based simulators to impose physics constraints, as well as massive databases of scientific data or knowledge bases. This approach has already proven successful in advancing scientific discovery, with applications spanning mechanical engineering \cite{ni_mechagents_2024, Buehler2023MechGPTModalities}, protein design \cite{ghafarollahi_protagents_2024, ghafarollahi2025sparks}, and materials design in areas such as bioinspired and inorganic systems \cite{luu_bioinspiredllm_2024,ghafarollahi2025sparks, luu_generative_2025, stewart_molecular_2025, stewart_higher-order_2026}.

Given these earlier results, there is growing interest in developing agent architectures that coordinate domain-specialist agents to address complex scientific challenges. Research directions such as group-chat structures, dynamic agent access, chat-memory, and adaptive task and tool allocation represent promising avenues for advancement.

An immediate priority is finding ways to extend agent access and use of large-scale knowledge sources. Open research opportunities remain for developing more efficient representations of knowledge in generative AI for science, as originally suggested in \cite{buehler_generative_2024,buehler_accelerating_2024}. A notable strategy is the incorporation of ontological knowledge graphs \cite{marquis_category_2023, mac1998categories, giesa_category_2012, markovic_optimizing_2025, sequeda_knowledge_2025,schilling-wilhelmi_text_2024,Buehler2025GraphPRefLexOR}, a structure in which concepts are represented as nodes and their relations as edges. Recent applications have demonstrated how LLMs can distill vast amounts of scientific text into these structured maps \cite{buehler_accelerating_2024,Buehler2025Chaos}. Graph  representations have also been utilized to develop graph-native reasoning models, as suggested in~\cite{Buehler2025GraphPRefLexOR}. Within agentic frameworks, such graphs become callable tools for in-context learning, enabling agents to retrieve knowledge dynamically and engage with massive scientific data \cite{ghafarollahi2025sciagents, stewart_higher-order_2026}. Because this approach does not rely on training, knowledge is maintained independently of model weights, ensuring transparency, facilitating updates, mitigating drift during inference, and reducing computational overhead. Knowledge graphs also allow retrieval-augmented generation (RAG) to be performed more efficiently than with unstructured text. Condensing large corpora into structured representations enables progress toward a scalable world model that integrates knowledge across domains. Further, relationships are represented explicitly in these graphs, making it possible to leverage the graph structure creatively to both uncover previously unknown connections and gain new insights into the organization of knowledge within a field.

A compelling application of this approach is the use of knowledge graphs as a substrate for idea mining \cite{buehler_generative_2024,buehler_accelerating_2024, ghafarollahi2025sciagents}. Analogous to mapping neurons and their synaptic connections, agents identify nodes within the graph and traverse the connectivities between them to uncover new relationships. This approach demonstrates the synergy of ontological knowledge graphs with the generative and reasoning capabilities of LLMs for idea generation. However, while single agents can contribute creative initial ideas, not every idea constitutes a robust hypothesis. 
%
Existing KG-based idea-mining approaches can surface novel associations, but they do not, by themselves, provide a reliable mechanism for (i) decomposing a design problem into verifiable sub-questions, (ii) grounding each sub-question in traceable evidence, and (iii) integrating cross-subtask constraints into a single testable hypothesis. This motivates an agentic architecture in which knowledge graphs are not only retrieval indices but controllable reasoning substrates.
%
To address this need, we introduce GraphAgents, a multi-agent framework that couples evidence-grounded retrieval with controllable graph traversal to support structured, traceable hypothesis generation and refinement for materials design.

\subsection{Model Framework}

Our framework integrates knowledge graph construction with a multi-agent system to support systematic hypothesis generation. Large corpora of scientific literature are distilled into two complementary knowledge graphs: a PFAS-Specific Knowledge Graph, for depth in the target domain, and a broader Material Properties Knowledge Graph for cross-domain exploration. These structured resources are coupled with a sequence of specialized agents that decompose design problems, retrieve and synthesize evidence from the collected corpus and developed knowledge graphs, extract key design parameters, and traverse graphs with heuristic algorithms to uncover novel relationships. The final stage integrates these insights into hypotheses that are both creative and grounded in prior knowledge. This combination of structured knowledge representation and coordinated agent reasoning provides a scalable foundation for scientific discovery. An overview of the full pipeline can be seen in Figure~\ref{fig:pipeline} with a profile of each agent shown in Figure~\ref{fig:agent_profiles}. 

This systematic approach is designed to emulate deep scientific thinking to accelerate the search for sustainable material alternatives. By combining domain expertise with methods that connect otherwise isolated pockets of knowledge, it becomes possible to broaden the search space, uncover overlooked candidate materials, and generate more creative, testable hypotheses.

\subsection{Outline of the Paper}

The subsequent sections of this paper are organized as follows. Section~\ref{sec:results} presents the results and discussion, including the design of the multi-agent system, knowledge graph construction, and hypothesis generation experiments using heuristic traversal algorithms. This section also reports ablation studies that evaluate the contributions of individual agents to the overall system performance. Section~\ref{sec:conclusion} concludes the paper with a summary of findings and discusses future directions for extending the framework. Finally, Section~\ref{sec:methods} provides details of the materials and methods, including model implementation, knowledge graph generation, and experimental setup.

\begin{figure}[H]
    \centering
    \includegraphics[width=1\linewidth]{figures/overview.pdf}
    \caption{Overview of the proposed agentic pipeline for materials discovery. The system begins with user input, which the Planner decomposes into sub-questions focusing on design parameters like material properties, operational conditions, failure modes, and interface requirements. The Hybrid GraphWeave agent retrieves evidence to answer each sub-question, from which the Evaluator selects pertinent design keywords with associated metrics. The Creative GraphWeave agent explores graph traversal pathways using the Evaluator keywords as stopping points to uncover novel connections and ideas. These are then synthesized by an Engineer into hypotheses that integrate both creative insights and contextual evidence. Bottom arrows illustrate how agents selectively retrieve and incorporate earlier responses as context for their own reasoning.}
    \captionsetup{skip=5pt}
    \label{fig:pipeline}
\end{figure}